%%文件: scaller.tex
%%作者: thegreatchaos@thegreatchaos.com
%%日期: 2026-08-26 15:07:31
通过向特定版本的\href{https://github.com/intel-innersource/applications.ai.gpu.vllm-xpu}{vLLM}, \href{https://github.com/vllm-project/vllm-xpu-kernels}{vllm-xpu-kernels}打patch + customized ESIMD kernels方式来enable features/optimization. 
相应的patches及esimd的kernels放在\href{https://github.com/intel/llm-scaler/}{llm-scaller}

开发流程:
\begin{enumerate}
    \item 从\href{https://github.com/intel-innersource/applications.ai.gpu.vllm-xpu}{vLLM}, \href{https://github.com/vllm-project/vllm-xpu-kernels}{vllm-xpu-kernels}, \href{https://github.com/intel/llm-scaler/}{llm-scaler}fork出三个repo: \href{https://github.com/intel-sandbox/llm-scaler-vllm-xpu}{vBase}, 
	\href{https://github.com/analytics-zoo/vllm-xpu-kernels}{kBase}, \href{https://github.com/intel/llm-scaler/}{sBase}
    \item 在vBase, kBase上修改, 如果涉及esimd kenrels则需要再sBase/vllm/custom-esimd-kernels-vllm目录下修改
    \item 完成开发/测试之后将vBase, kBase的改动以patch形式放到llm-scaler下, 如果有esimd kernels改动则merge到llm-scaler下
\end{enumerate}


既有xpu-kernels, esimd存在的必要?  xpu-kernels及sycl-tla过于heavy: 设计复杂, 各个算子独立. ESIMD是一个相对轻量的shortcut, 易于实现kernel fusion.

替代与fallback. 两者都是通过PyTorch的算子的方式注册, 被vLLM使用.
\begin{longtable}{p{9cm}|p{9cm}}
xpu-kernels&esimd\\
\hline
走vLLM官方的dispatch机制, 在vLLM的platform初始化时import vllm\_xpu\_kernels.\_C/\_moe\_C/\_xpu\_C. //注册rms\_norm/flash\_attn/grouped\_gemm...等算子
&
走显示注册机制, 在vLLM layer中显式import custom\_esimd\_kernels\_vllm(失败则用vllm中默认的实现), 同时也会读取环境变量:DISABLE\_ESIMD\_{GEMV/NORM/GEMM/...}来启用/禁用ESIMD的kernels\\
\end{longtable}


如何添加一个新的ESIMD算子, 以ESIMD\_GEMM算子为例
\begin{enumerate}
    \item vllm入口的修改. 显式替代xpu-kernels. git diff > vllm.patch
\begin{lstlisting}[language=python]
#vBase/vllm/model_executor/layers/esimd_utils.py中添加
_DISABLE_ESIMD_GEMM = _env_disabled("DISABLE_ESIMD_GEMM")
def disable_esimd_gemm() -> bool:
    return _DSIABLE_ESIMD_GEMM
#vBase/vllm/model_executor/kernels/linear/scaled_mm/xpu.py
from vllm.model_executor.layers.esimd_utils import(
    ...
    disable_esimd_gemm, #添加此行
    ...
)
#vBase/vllm/model_executor/kernels/linear/scaled_mm/xpu.py, 低精度矩阵乘法
class XPUFP8ScaledMMLinearKernel(FP8ScaledMMLinearKernel):
    ...
    def apply_weight(...):
	#添加以下
	if(ESIMD_AVAIABLE and and not disable_esimd_gemm()...):
	    esimd_gemm_fp8_pert(..., out)
	    return out;
	#添加结束
	return torch.ops._xpu_C.fp8_gemm_w8a16(...) #禁用/ESIMD not available则走xpu-kernels那套
#extra: 检查其他会用到fp8_gemm的部分, 视情况决定是否替换
#修改单元测试及benchmark.
# 1. grep -rn fp8_gemm_w8a16 -> 得到tests/kernels/quantization/test_xpu_fp8_scaled_mm.py, benchmarks/kernels/benchmark_xpu_fp8_w8a16.py, ... 修改这两个文件
\end{lstlisting}
    \item 目录结构, 文件如果不存在则创建他
\begin{lstlisting}
sBase/vllm/custom-esimd-kernels-vllm
|_csrc/xpu/esimd_kernel_gemm.sycl		//esimd实现文件
|_csrc/xpu/torch_extension_gemm.cc		//PyTorch C++算子注册, 之后可以在python中通过torch.ops来调用
|_include/kernel_ops.h				//esimd_gemm_fp8_pert的函数声明
|_python/custom_esimd_kernels_vllm/__init__.py	//让python初始化该package时候自动加载esimd_gemm_fp8_pert
|_python/custom_esimd_kernels_vllm/ops.py	//定义接口, 接口中通过torch.ops来调用C++的实现
\end{lstlisting}
    \item 'Makefile': setup.py
\begin{lstlisting}[language=c++]
ext_modules.append(  ####添加一个SyclExtension, 必须放在首个ext_modules = [...]之后
    SyclExtension(
        name="custom_esimd_kernels_vllm.custom_esimd_kernels_gemm",
        sources=[
            "csrc/xpu/esimd_kernel_esimd_kernel_gemm.sycl",
            "csrc/xpu/torch_extension_gemm.cc",
        ],
        include_dirs=[
            root / "include",
            root / "csrc",
        ],
        extra_compile_args={
            "cxx": ["-O3", "-std=c++17"],
            "sycl": ["-fsycl", "-ffast-math", "-fsycl-device-code-split=per_kernel",
                     "-fsycl-targets=spir64_gen", "-Xs", "-device bmg,ptl", ##这里添加目标平台AoT
                     f"-I{torch_include}"],
        },
        extra_link_args=["-Wl,-rpath,$ORIGIN/../../torch/lib"],
        py_limited_api=False,
    )
)
\end{lstlisting}
    \item 编译
\begin{lstlisting}
source /opt/intel/oneapi/setvars.sh
cd /path/to/scalerDev/ws/sBase/vllm/custom-esimd-kernels-vllm
TORCH_XPU_ARCH_LIST=ptl MAX_JOBS=1 python setup.py bdist_wheel
pip install --force--reinstall --no-deps dist/*.whl
\end{lstlisting}
    \item 测试
\begin{lstlisting}[language=python]
#scalerDev/ws/sBase/vllm/custom-esimd-kernels-vllm/tests/**.py #不同于上述的vllm中的测试, 这个实在custom-esimd-kernels-vllm中自己的UT
import sys, torch
import custom_esimd_kernels_vllm

def py_ref(...):
    return ..
def run(...):
    out = torch.ops.custom_esimd_kernels_vllm.esimd_newKernelName(...);
    diff = (out - py_ref(..))
    ok == diff < criteria_or_whatever
    if ok:
	print('pass')
    else:
	print('fail')

if __name__ == '__main__':
    cases = ...
    for c in case:
	run(c...);
\end{lstlisting}
\end{enumerate}

在esimd版本上的profiling, 从grouped\_gemm开始.
